% Workshop debugging_with_ai - ไฟล์สำหรับ Part B (มีจุดผิด 6 จุดโดยตั้งใจ)
% คอมไพล์ด้วย  latexmk -pdf main.tex  แล้วแก้ทีละจุดตามลำดับที่ข้อความผิดพลาดปรากฏ
% แก้ในบรรทัดเดิม อย่าเพิ่มหรือลบบรรทัด (หมายเลขบรรทัดจะตรงกับหน้า Workshop) และไม่ต้องแก้ refs.bib
\documentclass{article}
\usepackage[T1]{fontenc}
\usepackage{amsmaths}
\usepackage{amssymb}
\usepackage[backend=biber, style=authoryear]{biblatex}
\addbibresource{refs.bib}
\setlength{\emergencystretch}{1em} % lets long words in the bibliography fit
\title{The Birthday Problem}
\author{Typesetting Club}
\date{November 2026}

\begin{document}

\maketitle

\section{The Question}
How many people must be in a room before it is more likely than not
that two of them share a birthday? Most people guess a number close to
half of 365, but the answer is only 23. The question goes back to
\textcite{mises1939} and is now a standard example in probability
textbooks \parencite{feller1968}. We computed every number in this note
with a short script called birthday_table.py.

\section{Computing the Probability}
Ignore leap years and assume that all 365 birthdays are equally likely.
Let $p_n$ be the probability that $n$ people all have different
birthdays. The second person must avoid one day, the third must avoid
two days, and so on, so
\begin{align}
  p_n &= \frac{365}{365}\cdot\frac{364}{365}\cdots\frac{365-n+1}{365}
         \label{eq:product}\\
      &= \prod_{k=0}^{n-1}\Bigl(1-\frac{k}{365}\Bigr).
\end{equation}
Since $1-x \le e^{-x}$ for every $x \in \R$, equation~\eqref{eq:product}
gives the estimate
\begin{equation}
  p_n \le \exp\Bigl(-\frac{n(n-1)}{730}\Bigr).
  \label{eq:estimate}
\end{equation}
The right-hand side falls below $\frac{1}{2}$ when $n \ge 23$.

\section{Results}
Table~\ref{tab:value} lists the probability $1-p_n$ that at least two
of $n$ people share a birthday. With 70 people a shared birthday is
almost certain.

\begin{table}[ht]
  \centering
  \caption{Probability of a shared birthday}
  \label{tab:values}
  \begin{tabular}{rc}
    \hline
    $n$ & $1-p_n$ \\
    \hline
    10 & 0.117 \\
    20 & 0.411 \\
    23 & 0.507 \\
    30 & 0.706 \\
    50 & 0.970 \\
    70 & 0.999 \\
    \hline
  \end{tabular}
\end{table}

The result surprises us because we compare ourselves with the others,
while the problem counts every pair of people. With 23 people there are
already 253 pairs \parencite{diaconis1989}.

\printbibliography

\section*{Verification Log}
% Part B: เติมแถว B2-B6 ของตารางนี้ (แถว B1 เป็นตัวอย่าง)
\begin{tabular}{|l|p{4cm}|p{1.6cm}|p{4cm}|}
  \hline
  \textbf{Bug} & \textbf{What the AI said} & \textbf{Correct?} & \textbf{What I did} \\
  \hline
  B1 & The package name is misspelled. & Yes & Wrote amsmath on line 6. \\
  \hline
  B2 & & & \\
  \hline
  B3 & & & \\
  \hline
  B4 & & & \\
  \hline
  B5 & & & \\
  \hline
  B6 & & & \\
  \hline
\end{tabular}

\end{document}
